\section{Critical comparison and recovery}
\label{sec:uniqueness-factor}

We now show that two conditional copies with the same innovations agree
at the critical parameter $\theta=b^{-1/2}$, where $b=d-1\ge2$.
The weighted path estimates reduce the question to the variance
$V(t)=\E[(\Delta h_o(t))^2]$ of the site-field difference.
Their principal noise term has the scale
\[
 S(t)\sum_{r\ge1}r e^{-c\sqrt t\,r}\asymp \frac{S(t)}{t}
 \qquad(t\downarrow0),
\]
where $S(t)=\E[D_e(t)^2]$ and
$D_{(i,j)}=\log a(h_{j\to i})-\log a(h'_{j\to i})$.
The earlier estimate $S\le CV$ would leave a coefficient $1/t$,
which is not integrable at zero.  We improve it by using the fact that
$(\log a)'$ vanishes at the origin.  A field is usually small at early
times; the first result below makes the exceptional probability small
enough to compare it with any positive value of $V$.

\subsection{All moments and a small-field tail}

We use the fixed-time posterior identities and critical information
estimate already proved for one observation process.  At a deterministic
time $t>0$, fix an oriented edge $v\to w$ and write
\[
 Y=h_{v\to w}(t),\qquad M(t)=\E\tanh^2Y.
\]
The law of $Y$ is independent of the choice of edge and is symmetric.
Recall that $M(t)\le C\sqrt t$ for $0<t\le1$, that
$|Y-X_v(t)|\le b\beta$, and that
$X_v(t)=t\sigma_v+\sqrt t\,G$ with $G$ a standard normal variable
independent of the spin.  These are statements about the planted
observation law.  Each marginal of the conditional coupling has this
same law.

The second moment is critical because $b\theta^2=1$, whereas
$b\theta^p<1$ for every integer $p\ge3$; this strict contraction
lets the recursion propagate the second-moment scale to every higher
integer order.

\begin{lemma}[Critical field moments and tail]
\label{lem:critical-tail}
There are constants $0<K<\infty$ and $t_0\in(0,1)$, depending only on
$d$, such that for every integer $p\ge1$ and every $0<t\le t_0$,
\begin{equation}
 \E|h_{v\to w}(t)|^p\le e^{Kp^2}t^{p/4}.
 \label{eq:critical-all-moments}
\end{equation}
There are also $c>0$ and $t_{\mathrm{tail}}\in(0,t_0]$ such that
\begin{equation}
 \P\bigl(|h_{v\to w}(t)|>t^{1/8}\bigr)
 \le \exp\bigl(-c\log^2(1/t)\bigr),
 \qquad 0<t\le t_{\mathrm{tail}}.
 \label{eq:critical-field-tail}
\end{equation}
Both inequalities hold as well when conditioned on either value of the
component-root spin $\sigma_v$, or on either value of the deleted-neighbor
spin $\sigma_w$.
\end{lemma}

\begin{proof}
We begin with the second moment.  On $\{|X_v(t)|\le1\}$ the cavity
bound gives $|Y|\le1+b\beta$.  Comparison on this compact interval
therefore gives $Y^2\le C\tanh^2Y$.  On its complement we use
$Y^2\le2X_v(t)^2+2b^2\beta^2$.  For $t\le1/2$, the event
$|t\sigma_v+\sqrt t\,G|>1$ requires $|G|>1/(2\sqrt t)$.
Integration of the normal density, decreasing the positive exponent
constant if necessary, yields
\[
 \E\bigl[(1+X_v(t)^2)\mathbf 1_{\{|X_v(t)|>1\}}\bigr]
 \le C e^{-c/t}.
\]
Combining the two events with $M(t)\le C\sqrt t$ proves that, after
fixing some $t_0\le1/2$,
\begin{equation}
 A_2(t):=\E|Y|^2\le C\sqrt t,
 \qquad A_1(t):=\E|Y|\le C^{1/2}t^{1/4}
 \quad(0<t\le t_0).
 \label{eq:critical-moment-base}
\end{equation}
The second inequality is Cauchy--Schwarz.  The cavity bound also shows
that $A_p(t):=\E|Y|^p$ is finite for every integer $p\ge1$.

We next explain why the recursion can be expanded using independent
terms without changing these absolute moments.  For an integrable
function $f$, the spin-tilt identities from the posterior calculation
are
\begin{align*}
 \E[f(Y)\mid\sigma_v=\eta]
   &=\E[f(Y)(1+\eta\tanh Y)],\\
 \E[f(Y)\mid\sigma_w=\eta]
   &=\E[f(Y)(1+\eta\theta\tanh Y)],
 \qquad \eta\in\{-1,1\}.
\end{align*}
These are the identities \eqref{eq:component-field-tilts} from
Lemma~\ref{lem:component-spin-tilts}.
Because $Y$ has a symmetric marginal law, the term containing
$\tanh Y$ integrates to zero against every even function of $Y$.
In particular conditioning on either spin leaves each $A_p(t)$
unchanged.

Condition now on $\sigma_v=+1$ in the cavity recursion at $v$.
The own observation and the $b$ child-component observations are
independent under this conditional law: their spin transitions and
their planted noises lie in disjoint branches given the root spin.
The own observation has law $t+\sqrt t\,G$.  A child field is
conditioned through its deleted-neighbor spin, so the even-function
consequence of the second tilt gives, for every integer $r\ge1$,
\begin{equation}
 \E\bigl[|\psi(h_{u\to v}(t))|^r\mid\sigma_v=+1\bigr]
   =\E|\psi(Y)|^r\le\theta^r A_r(t).
 \label{eq:child-absolute-moments}
\end{equation}
Here $\psi(0)=0$ and $0<\psi'(z)\le\theta$ imply
$|\psi(z)|\le\theta|z|$.

There is an absolute constant $K_0$ such that
\begin{equation}
 \E|t+\sqrt t\,G|^r\le e^{K_0r^2}t^{r/4}
 \qquad(r\ge1\text{ an integer},\ 0<t\le1).
 \label{eq:own-gaussian-moments}
\end{equation}
Indeed, $|t+\sqrt t\,G|\le\sqrt t(1+|G|)$ in this time range,
and $\E|G|^r\le(C\sqrt r)^r$.  The inequality
$(x+y)^r\le2^{r-1}(x^r+y^r)$ for nonnegative $x,y$ bounds the
remaining constant by $e^{K_0r^2}$, uniformly over integer $r$.
We also used $t^{r/2}\le t^{r/4}$.

Choose $K\ge2K_0$ large enough that \eqref{eq:critical-all-moments}
holds for $p=1,2$, using \eqref{eq:critical-moment-base}.  Suppose it
holds for all positive integers below $p$, where $p\ge3$.
The triangle inequality in the cavity recursion, followed by the
multinomial expansion and the conditional independence just established,
bounds $A_p(t)$ by the expectation of the $p$th power of the sum of
the absolute values of the own term and the $b$ child terms.
The terms with exponent $p$ on a single child contribute at most
$b\theta^p A_p(t)$.  The term supported on the own observation is at
most $e^{Kp^2/2}t^{p/4}$ by \eqref{eq:own-gaussian-moments}.

Every other exponent vector $(r_0,r_1,\ldots,r_b)$ has sum $p$ and
at least two positive entries.  Hence every positive entry is below
$p$, and the induction hypothesis applies to its child factors.
Using \eqref{eq:own-gaussian-moments} for its own factor, if present,
and discarding powers of $\theta\le1$, its product of moments is at
most
\[
 e^{K\sum_{i=0}^b r_i^2}t^{p/4}.
\]
For such a vector,
$\sum_i r_i^2\le(p-1)^2+1=p^2-2(p-1)$: merging positive entries
until only two remain, then moving units to the larger one, can only
increase the sum of squares.  The sum of all multinomial coefficients
in $b+1$ slots is $(b+1)^p$.  Thus
\begin{equation}
 (1-b\theta^p)A_p(t)
 \le e^{Kp^2}t^{p/4}
   \left[e^{-Kp^2/2}+(b+1)^p e^{-2K(p-1)}\right].
 \label{eq:critical-moment-induction}
\end{equation}
The coefficient on the left is positive uniformly over $p\ge3$:
\[
 1-b\theta^p=1-b^{1-p/2}\ge\kappa_b:=1-b^{-1/2}>0.
\]
Enlarge the same constant $K$ once more so that
\[
 2K>\log(b+1),\qquad
 e^{-9K/2}+(b+1)^3e^{-4K}\le\kappa_b.
\]
Both terms in brackets in \eqref{eq:critical-moment-induction}
decrease as the integer $p$ increases from $3$.  Their sum is
therefore at most $\kappa_b\le1-b\theta^p$ for every such $p$.
Absorption proves \eqref{eq:critical-all-moments} with this single
$K$, for every integer order and all $t\le t_0$.

To obtain the tail, put $L=\log(1/t)$ and choose the integer
$p=\lfloor L/(16K)\rfloor$.  If $L\ge32K$, then
$L/(32K)\le p\le L/(16K)$.  Markov's inequality gives
\[
 \P(|Y|>t^{1/8})
 \le e^{Kp^2-pL/8}
 \le e^{-pL/16}
 \le e^{-L^2/(512K)}.
\]
Take $t_{\mathrm{tail}}\le\min(t_0,e^{-32K})$.
Both $|Y|^p$ and the indicator of the tail event are even functions
of $Y$, so the spin-tilt identities prove all the conditional
assertions as well.
\end{proof}

\subsection{Comparing log-factor and site-field differences}

Fix a finite horizon $T>0$.  The conditional-copy construction gives
\[
 |\Delta h_v(t)|\le H_T:=2T+2d\beta,
 \qquad 0\le V(t)\le H_T^2\quad(0\le t\le T),
\]
and Lemma~\ref{lem:copy-variance} makes $V$ absolutely continuous,
with $V(0)=0$.  Choose a constant, fixed throughout this horizon,
\begin{equation}
 K_T>\max\{1,H_T^2\},\qquad
 I(t)=1+\log\frac{K_T}{V(t)}\quad\text{when }V(t)>0.
 \label{eq:osgood-size}
\end{equation}
In particular $I(t)>1$, including at times when $V(t)\ge1$.

\begin{lemma}[Small-time log-factor estimate]
\label{lem:small-log-loss}
There is a fixed $t_*\in(0,\min\{t_{\mathrm{tail}},e^{-1}\}]$ such
that, for $0<t\le\min\{t_*,T\}$ with $V(t)>0$,
\begin{equation}
 S(t)\le C_T V(t)
   \left(t^{1/4}+\frac{I(t)}{\log^2(1/t)}\right).
 \label{eq:small-log-loss}
\end{equation}
At every time in $[0,T]$, the bound $S(t)\le CV(t)$ remains valid.
\end{lemma}

\begin{proof}
Write $Y=h_{j\to i}(t)$ and $Y'=h'_{j\to i}(t)$ for the edge used to
define $S(t)$.  The two cavity equations on this edge and
$|\psi'|\le\theta$ give the earlier comparison
\[
 |Y-Y'|\le\frac{|\Delta h_j|+\theta|\Delta h_i|}{1-\theta^2}.
\]
Invariance and Cauchy--Schwarz consequently give
\begin{equation}
 \E|Y-Y'|^2\le\frac{V(t)}{(1-\theta)^2},
 \qquad |Y-Y'|\le\frac{H_T}{1-\theta}.
 \label{eq:critical-component-comparison}
\end{equation}
Also
\[
 |(\log a)'(z)|
 =\frac{2(1-\theta^2)|\tanh z|}{1-\theta^2\tanh^2z}
 \le2|\tanh z|\le2.
\]
The global bound gives $S(t)\le4V(t)/(1-\theta)^2$ and
$|D_e(t)|\le2H_T/(1-\theta)$.

On the event $|Y|,|Y'|\le t^{1/8}$ the segment between them also lies
in that interval, and $|(\log a)'(z)|\le2t^{1/8}$ there.
The mean-value theorem and \eqref{eq:critical-component-comparison}
therefore give
\[
 \E[D_e(t)^2;\ |Y|,|Y'|\le t^{1/8}]
 \le\frac{4}{(1-\theta)^2}t^{1/4}V(t).
\]
For the complementary event the global square-integrable bound gives
$CV(t)$, while the deterministic bound on $D_e$ and a union bound
using \eqref{eq:critical-field-tail} in the two marginals give
$C_Te^{-c\log^2(1/t)}$.  Thus
\begin{equation}
 \E[D_e(t)^2;\ \max(|Y|,|Y'|)>t^{1/8}]
 \le\min\{CV(t),C_Te^{-c\log^2(1/t)}\}.
 \label{eq:bad-field-minimum}
\end{equation}
The argument uses only the two marginal probabilities; it requires
no independence between the two fields.

Set $u=V(t)>0$ and $J=\log^2(1/t)\ge1$.  If $u\le e^{-cJ/2}$,
then $I(t)\ge cJ/2$, so the first member of the minimum in
\eqref{eq:bad-field-minimum} is at most $C u I(t)/J$.
If $u>e^{-cJ/2}$, the second member divided by $u$ is at most
$C_Te^{-cJ/2}$.  Since $J e^{-cJ/2}$ is bounded for $J\ge1$
and $I(t)\ge1$, this is at most $C_T I(t)/J$.
Adding the small-field contribution proves \eqref{eq:small-log-loss}.
\end{proof}

\subsection{A cutoff depending on the discrepancy}

We next sum the weighted estimates.  Recall
$\omega(t)=\min\{\sqrt t,1\}$, and suppress $t$ at a fixed time
with $0<t\le T$ and $V(t)>0$.  The preceding path results give
\begin{align*}
 Q_r&\le C rS,\qquad L_r\le V &&(r\ge1),\\
 Q_r&\le C e^{-c\omega r}rS+C_T r^2e^{-\alpha b^{\gamma r}}
 &&(r\ge C_0'/\omega),\\
 L_r&\le C e^{-c\omega r}V+C_T e^{-\alpha b^{\gamma r}}
 &&(r\ge C_0'/\omega),
\end{align*}
for fixed positive constants $c,\alpha,\gamma,C_0'$ depending only on
$d$.  The last two lines follow from Proposition~\ref{prop:weighted-paths}, with
constants decreased or enlarged so that both exceptional terms have
the indicated common form.  These estimates include weights and
discrepancies formed from the same observations.

Choose $C_1$ sufficiently large that
\[
 C_1\ge C_0',\qquad \gamma C_1\log b\ge1,
 \qquad (\alpha/2)b^{\gamma C_1}\ge2,
\]
and take the integer
\begin{equation}
 R=\left\lceil C_1\bigl(\omega^{-1}+\log I\bigr)\right\rceil.
 \label{eq:variance-cutoff}
\end{equation}
The weighted estimates apply for $r>R$.  Since $\omega^{-1}\ge1$
and $I>1$, the choices above also give
\[
 (\alpha/2)b^{\gamma r}
 \ge (\alpha/2)b^{\gamma C_1\omega^{-1}}
                  I^{\gamma C_1\log b}
 \ge 2I \qquad(r>R).
\]
Splitting the exceptional exponent into equal halves therefore gives
\begin{equation}
 e^{-\alpha b^{\gamma r}}
 \le e^{-2I}e^{-(\alpha/2)b^{\gamma r}}
 =e^{-2}\left(\frac{V}{K_T}\right)^2
       e^{-(\alpha/2)b^{\gamma r}}
 \le V^2 e^{-(\alpha/2)b^{\gamma r}}\quad(r>R).
 \label{eq:exception-absorption}
\end{equation}
The last step uses $K_T>1$, and so remains valid when $V\ge1$.

\begin{lemma}[The noise and drift path sums]
\label{lem:cutoff-sums}
At each time $0<t\le T$ with $V(t)>0$,
\begin{equation}
 \begin{split}
 \sum_{r\ge1}Q_r
 &\le C_T\left[S\bigl(\omega^{-2}+(\log I)^2\bigr)+V^2\right],\\
 \sum_{r\ge1}\sqrt{Q_r}
 &\le C_T\left[\sqrt S\bigl(\omega^{-3/2}
                          +(\log I)^{3/2}\bigr)+V\right],\\
 \sum_{r\ge1}L_r
 &\le C_T\left[V\bigl(\omega^{-1}+\log I\bigr)+V^2\right].
 \end{split}
 \label{eq:cutoff-sums}
\end{equation}
\end{lemma}

\begin{proof}
For $r\le R$, the estimates $Q_r\le CrS$ and $L_r\le V$ give
the respective bounds
\[
 CS\sum_{r\le R}r\le CSR^2,
 \qquad C\sqrt S\sum_{r\le R}\sqrt r\le C\sqrt S R^{3/2},
 \qquad \sum_{r\le R}L_r\le VR.
\]
For the ordinary exponential terms at $r>R$, use
\begin{equation}
 \sum_{r\ge1}r^p e^{-c\omega r}
 \le C_{p,c}\omega^{-p-1}
 \quad\left(p\in\left\{0,\tfrac12,1\right\},\ 0<\omega\le1\right).
 \label{eq:exponential-length-sums}
\end{equation}
For example, compare the sum with
$\int_0^\infty(x+1)^p e^{-c\omega x}\,\mathrm dx$ and put
$y=\omega x$; since $\omega\le1$, the resulting integral is bounded
by $\omega^{-p-1}\int_0^\infty(y+1)^p e^{-cy}\,\mathrm dy$.
This gives $CS\omega^{-2}$ for the first sum and $CV\omega^{-1}$
for the third.  For the middle sum use
$\sqrt{x+y}\le\sqrt x+\sqrt y$; the ordinary exponential contribution
is bounded by $C\sqrt S\sum_{r\ge1}r^{1/2}e^{-c\omega r/2}
\le C\sqrt S\omega^{-3/2}$.

For the exceptional terms, \eqref{eq:exception-absorption} bounds
their three contributions by constants times
\[
 V^2\sum_{r\ge1}r^2e^{-(\alpha/2)b^{\gamma r}},\qquad
 V\sum_{r\ge1}r e^{-(\alpha/4)b^{\gamma r}},\qquad
 V^2\sum_{r\ge1}e^{-(\alpha/2)b^{\gamma r}},
\]
respectively.  Each series is finite because its exponential exponent
grows exponentially in $r$.  In the middle series the square root of
the $V^2$ factor is $V$, which is why the drift sum has this order.
Finally $\omega^{-1}+\log I\ge1$, so the ceiling in
\eqref{eq:variance-cutoff} gives
$R\le(C_1+1)(\omega^{-1}+\log I)$.
Bounding each power of this sum by a constant times the sum of its
corresponding powers proves \eqref{eq:cutoff-sums}.
\end{proof}

\subsection{An integrable comparison at time zero}

\begin{proposition}[Critical conditional-copy uniqueness]
\label{prop:critical-uniqueness}
For the invariant conditional coupling with common innovations at
$\theta=b^{-1/2}$, the site-field discrepancy satisfies
$V(t)=0$ for every $t\ge0$.
\end{proposition}

\begin{proof}
Work on the fixed horizon $[0,T]$ used above.  Inserting the noise
and drift bounds \eqref{eq:noise-bound} and~\eqref{eq:drift-bound}
into the variance identity~\eqref{eq:variance}, and using
$n_r\theta^{2r}=db^{r-1}b^{-r}=d/b$ for $r\ge1$, gives
\begin{equation}
 V'(t)\le C\left(
 V(t)+\sum_{r\ge1}Q_r(t)+\sum_{r\ge1}L_r(t)
            +\sqrt{V(t)}\sum_{r\ge1}\sqrt{Q_r(t)}\right)
 \label{eq:variance-path-reduction}
\end{equation}
for almost every $t$.  The sums in Lemma~\ref{lem:cutoff-sums} are
finite at the positive times with $V(t)>0$ where we now use this
inequality.

Insert those three bounds into \eqref{eq:variance-path-reduction}.
Use $S\le CV$, and hence $\sqrt{VS}\le C V$, for all terms except
$S\omega^{-2}$.  The bounds $V\le H_T^2$ imply
$V^{3/2}\le H_TV$ and $V^2\le H_T^2V$.  Finally
\[
 (\log I)^p\le C_p I\qquad(I\ge1,\ 0<p\le2),
\]
since $x^p e^{-x}$ is bounded on $x\ge0$.  As $\omega^{-1}\le
\omega^{-3/2}$, all the terms just treated are at most
\[
 C_TV\bigl(I+\omega^{-3/2}\bigr).
\]
This identifies the worst drift contribution explicitly:
$\sqrt{VS}\,\omega^{-3/2}\le CV\omega^{-3/2}$.

For $0<t\le\min\{t_*,T\}$, we have $\omega=\sqrt t$.
Lemma~\ref{lem:small-log-loss} gives the remaining estimate
\[
 S\omega^{-2}
 \le C_TV\left(t^{-3/4}
           +\frac{I}{t\log^2(1/t)}\right).
\]
Together with $I\ge1$, the preceding bounds prove
\begin{equation}
 V'(t)\le k_T(t)V(t)
        \left(1+\log\frac{K_T}{V(t)}\right)
 \quad\text{for almost every $t\in(0,T]$ with $V(t)>0$},
 \label{eq:integrable-osgood}
\end{equation}
where, after increasing $C_T$, we may take
\begin{equation}
 k_T(t)=
 \begin{cases}
 C_T\left(1+t^{-3/4}
          +\dfrac{1}{t\log^2(1/t)}\right),&0<t\le\min\{t_*,T\},\\[2mm]
 C_T,&t_*<t\le T.
 \end{cases}
 \label{eq:osgood-coefficient}
\end{equation}
For the second branch, $\omega(t)\ge\sqrt{t_*}>0$ because
$t_*\le1$; there we use $S\le CV$ also for $S\omega^{-2}$.
Thus that branch is bounded on the whole later interval.
If $T\le t_*$ there is no later interval.  The small-time expression
is used only for $t\le t_*\le e^{-1}$, and it is integrable there:
\[
 \int_0^{t_*}t^{-3/4}\,\mathrm dt=4t_*^{1/4},\qquad
 \int_0^{t_*}\frac{\mathrm dt}{t\log^2(1/t)}
 =\frac{1}{\log(1/t_*)}.
\]
In particular $k_T\in L^1(0,T)$.

We finish with an Osgood-type comparison~\cite{Osgood1898}, proving
the needed absolutely continuous form directly.  Suppose that $V$ is positive somewhere
on $[0,T]$.  Choose such a time $t$, and let $a<t$ be the left
endpoint of its component in the positivity set of $V$.
Continuity and $V(0)=0$ give $V(a)=0$.  On any closed interval
$[u,t]$ with $a<u<t$, continuity bounds $V$ away from zero.
The function $z\mapsto\log(1+\log(K_T/z))$ has bounded derivative
on the range of $V$ on that interval.  Its composition with the
absolutely continuous function $V$ is therefore absolutely continuous,
and the chain rule and \eqref{eq:integrable-osgood} give
\[
 (\log I)'=-\frac{V'}{VI}\ge-k_T\quad\text{almost everywhere on }[u,t].
\]
Integrating yields
\[
 \log I(u)\le\log I(t)+\int_u^t k_T(s)\,\mathrm ds
             \le\log I(t)+\int_a^t k_T(s)\,\mathrm ds<\infty.
\]
The last bound is independent of $u$.  But $V(u)\to V(a)=0$ as
$u\downarrow a$, so $I(u)\to\infty$ and $\log I(u)\to\infty$,
a contradiction.  Hence $V$ vanishes on $[0,T]$.  Since the horizon
was arbitrary, it vanishes at every time.
\end{proof}
